%% ============================================================
%%  Abstract
%% ============================================================
\chapter*{Abstract}
\addcontentsline{toc}{chapter}{Abstract}

Modern networked environments face an increasingly sophisticated threat landscape
in which adversaries pursue structured campaigns across multiple exploitation
phases.
Traditional intrusion-detection and honeypot management systems rely on static
rule bases or fixed-threshold anomaly detectors that cannot adapt as attacker
strategies evolve in real time.
This thesis presents \textbf{HoneyIQ}, a simulation framework and decision
policy for adaptive honeypot management grounded in the Lockheed~Martin
Cyber Kill Chain~\citep{hutchins2011intelligence}.

The attacker subsystem models adversarial behaviour as a pair of coupled
Discrete-Time Markov Chains governing attack-type and kill-chain-stage
transitions.
Network traffic features are generated by sampling from parametric
distributions calibrated against the UNSW-NB15 intrusion detection
benchmark~\citep{moustafa2015unsw}.
Four attacker intent profiles --- \emph{Stealthy}, \emph{Aggressive},
\emph{Targeted}, and \emph{Opportunistic} --- are encoded as intent-specific
transition modifiers, producing qualitatively distinct attack campaigns from
a shared base model.

The principal contribution of this thesis is the
\textbf{Stage-Escalation Decision Matrix (SEDM)}, a transparent, deterministic
policy that maps the current kill chain stage and a Markov-chain-derived
escalation risk score to one of five honeypot responses
(ALLOW, LOG, TROLL, BLOCK, ALERT).
A composite threat-level signal integrating attack severity, kill chain
progress, escalation rate, and cumulative attack pressure provides a
single normalised index for situational awareness.
A Random Forest classifier~\citep{breiman2001random} provides attack-type
identification from raw network features as a supporting component.

Experimental evaluation across 30~episodes per attacker intent (200 steps
each) demonstrates that the SEDM achieves consistently high threat coverage:
detection rates of \textbf{99.09\%}, \textbf{99.47\%}, \textbf{99.48\%},
and \textbf{99.41\%} for the Stealthy, Aggressive, Targeted, and
Opportunistic profiles respectively.
False positive rates range from 3.33\% (Targeted) to 35.56\% (Stealthy),
reflecting the inherent trade-off between sensitivity and specificity across
intent profiles with different threat-level distributions.
The SEDM's transparency, zero training overhead, and consistent cross-intent
performance are compared with the Deep Q-Network (DQN) baseline, whose
training on Opportunistic traffic achieves comparable detection rates
but lacks the interpretability required for operational deployment.

\vspace{1em}
\noindent\textbf{Keywords:} cyber kill chain, honeypot management, Stage-Escalation
Decision Matrix, deep reinforcement learning, Deep Q-Network, Markov chain,
intrusion detection, attacker intent, UNSW-NB15, adaptive cybersecurity.
\clearpage
